\section{Introduction}

% Prediction vs Causality
Economists are faced with two related, but very different type of questions. There is the question about prediction and there is the question about causality. This duality can be better illustrated by an example: the demand for hotel rooms. If we were given past data on room prices and occupancy of hotels, we would encounter a positive correlation between them: whenever the price is high, occupancy is expected to be high as well. Knowing nothing else, if you were asked to predict the price of a hotel whose occupancy is high, you should predict the price to be high as well. A naïve hotel owner, however, might conclude that by raising prices, he will increase the occupancy of his establishment. This, of course, doesn't make sense in the light of economic theory. The hotel owner is mixing up correlation with causation, prediction with causality. Demand for a good decreases as its price goes up. In the hotel case, there is a third factor that drives both demand and prices up. Vacation season increases demand for hotel rooms. Seeing this, hotel owners by raising prices. Thus, it becomes clear that the surge in demand due to high season is strong enough to make occupancy increase \textit{despite} the simultaneous price increase.

The search for causality has been long one of the main concerns of economists. In recent years, machine learning - a field that has been mostly concerned with prediction - started to be used together with econometrics to answer questions about causal inference. Following these new developments, we will reassess in this paper the question about the causal effect of Medicaid coverage on visits to the emergency department. This problem has been already analyzed with battle-tested econometric methods by \cite{taubmanMedicaidIncreasesEmergencyDepartment2014} . Here, however, we will provide an overview about causal inference, heterogeneitey in causal effects and how decision trees can be stiched together to uncover causal effects. Finally, we will apply causal forests, a method whose basic structure was developed within machine learning (\cite{breimanRandomForests2001}), to try to identify heterogeneity in the response to Medicaid coverage.


% Econometrics we will be analyzing a Machine learning and econometrics approach their fundamental questions they are trying to answer in different ways. At the heart of econometrics lies the constant search for an explanation. Econometricians won't think twice before sacrificing goodness of fit of their models to achieve a better, more precise estimate of the parameter whose ground truth they are trying to uncover. Machine learning focuses rather on prediction. It has developed during the last decades various techniques, heuristics and methods that have made it very successful in different fields.

\newpage

